\section{Conclusion}
\label{sec:conclusion}

% Written 29 Oct 2026 from the rc1 numbers.

The CIDR version of \duckpgq answered shortest-path queries with UDFs, and the UDFs inherited \duckdb's
tuple-driven parallelism, so a query with few pairs ran on one core and every query rebuilt the CSR. This paper
removed both limits without leaving the extension. The path-finding operator takes its parallelism from a CSR
partitioned by destination range and runs every traversal as a chain of events in \duckdb's scheduler. The CSR index
keeps the graph across queries and stays consistent with \duckdb's MVCC through epochs and a change log whose
visibility \duckdb itself decides, once per change group. Because the operator answers many pairs with one
traversal, the gap to the native graph systems grows with the number of pairs. At SF100 one query for 65\,536 pairs
takes 3.5\,s, against 155\,s for Neo4j and 169\,s for Kuzu, and at SF300 it takes 10\,s while every other system
exceeds ten minutes. Under the LDBC update stream a commit costs 15\,ms, against 342\,ms for Kuzu, and every
transaction sees exactly the edges it may see. \todo{Confirm the numbers against rc2.}

\duckdb's extension interfaces were enough, and where they fell short the gaps are specific. An index callback does
not say which transaction or commit made a change, and nothing can ask about visibility in the past. The scheduler has
no tasks for index work, so a fold runs in the transaction that triggers it. An extension operator can report a
cardinality but not a cost. Each gap cost us a mechanism in this paper, and each would be small to close in the
host system.

In the future, we will replace the thresholds of Section~\ref{sec:cost} with a cost model calibrated on these runs,
move the fold into the background, and extend the operator to weighted shortest paths, to an index over several edge
tables of one property graph, and to path reconstruction in the multi-source traversal.
